\section{实操步骤：如何运行 AutoResearch}

以下是运行 autoresearch 的完整 6 步流程：

\begin{enumerate}
    \item \textbf{Step 1：下载 skill}。获取 autoresearch skill，放进 Claude Code 或 Cowork 的 skills 文件夹。

    \item \textbf{Step 2：选一个要改进的 skill}。说``run autoresearch on my [skill name] skill''。选最让你头疼的那个——时好时坏、输出不稳定的 skill。

    \item \textbf{Step 3：Agent 问你 3 件事}：
    \begin{itemize}
        \item 要优化哪个 skill
        \item 用什么测试输入（例如``write landing page copy for an AI productivity tool''）
        \item 你的 checklist 问题是什么
    \end{itemize}

    \item \textbf{Step 4：跑一遍 skill，给出起始分数}。这是基准线（baseline）。作者的落地页 skill 起步只有 56\%——标题模糊、buzzword 泛滥、CTA 软弱，超过一半的检查项都没通过。

    \item \textbf{Step 5：浏览器弹出实时 dashboard}。包含以下内容：
    \begin{itemize}
        \item 分数随时间变化的曲线图
        \item 每项 checklist 的通过/失败明细
        \item 每次改动的日志
        \item 每 10 秒自动刷新
    \end{itemize}

    \item \textbf{Step 6：走开}。Agent 进入循环：
    \begin{itemize}
        \item 分析哪些检查项在失败
        \item 对 skill prompt 做一个小改动
        \item 重新测试
        \item 分数升就留下改动，分数降就撤销
        \item 一直自主运行，直到你叫停，或连续三次达到 95\%+ 为止
    \end{itemize}
\end{enumerate}

\begin{importantbox}{全自动运行，原始 Skill 不受影响}
你可以盯着 dashboard 看，也可以完全走开。Agent 不需要你参与任何决策。改进后的版本保存为单独的文件，你的\textbf{原始 skill 始终不会被修改}。
\end{importantbox}

\subsection{本章小结}

实操分为 6 步：下载 skill $\to$ 选目标 $\to$ 回答 3 个问题 $\to$ 建立基准线 $\to$ 观察 dashboard $\to$ 放手让 agent 跑。整个过程除了初始配置外完全自动化，且原始 skill 始终安全。

